\chapter{Istruzioni d'uso}

\section{Build dei test}
Per eseguire i test unitari:
\begin{lstlisting}[style=cppstyle]
make test
\end{lstlisting}
Questo comando invoca lo script e produce l'eseguibile di test nella cartella \texttt{test/build}.

Per pulire:
\begin{lstlisting}[style=cppstyle]
make test-clean
\end{lstlisting}
La pulizia e consigliata quando si cambiano toolchain o si aggiornano le dipendenze.

\section{Compilazione Arduino}
Per compilare lo sketch tramite lo script:
\begin{lstlisting}[style=cppstyle]
./test/run_tests.sh --with-arduino --board arduino:avr:uno
\end{lstlisting}
Questa modalita esegue prima i test unitari e poi compila lo sketch, riducendo il rischio di caricare un firmware non verificato.

Per caricare sulla board reale:
\begin{lstlisting}[style=cppstyle]
./test/run_tests.sh --upload --board arduino:avr:uno --port /dev/ttyACM0
\end{lstlisting}
Prima dell'upload, verificare che la porta seriale sia corretta.

\section{Simulazione Wokwi}
Il cablaggio di simulazione e definito in \texttt{diagram.json}. Aprire il progetto in Wokwi e avviare la simulazione per vedere il comportamento complessivo. La simulazione e utile per validare il flusso senza hardware fisico.
